% =====================================================================
%  Chapter 17b — Pressure-testing the landscape (S53–S57)
% =====================================================================
\chapter{Pressure-Testing the Landscape}
\label{ch:pressure-testing}

\begin{record}
\textbf{Sources of record.} Five session logs, read in full, all in
\texttt{SESSION-LOG/}:

\smallskip
{\raggedright\ttfamily
2026-09-12-S53-typeII-reconnaissance.md\\
2026-09-12-S54-wpt2-openai-profile.md\\
2026-09-12-S55-wpt3-forcing-floor.md\\
2026-09-12-S56-strategic-audit-literature.md\\
2026-09-13-S57-wpa-grujic-bmo-pricing.md\par}

\smallskip
Each is cross-checked here against its
\texttt{PROGRESS.md} \textsc{key finding} block
and, where finer detail was needed, against the proof document directly:
\S\,\texttt{sec:typeII-wpt1} (\Sn{53}), \S\,\texttt{sec:wpt2-openai}
(\Sn{54}), \S\,\texttt{sec:wpt3-forcing-floor} (\Sn{55}) and
\S\,\texttt{sec:wpa-s57} (\Sn{57}).

\smallskip
\textbf{Citation convention, as everywhere in this book.} Objects of the proof
document are named as text --- \texttt{lem:wpt1-nueff} and the rest --- because
that document is not this one and its labels do not resolve here. Only
cross-references \emph{inside} this book are set as references.

\smallskip
\textbf{No numerics.} This entire thread produced no experiment and no
\texttt{results.db} row. Everything in it is derivation and reading, and every
quantitative statement below traces to a session log or to a labelled object in
the proof document.
\end{record}

Chapter~\ref{ch:openai} surveyed the external landscape as it stood at \Sn{49}:
a forced blow-up construction, two strands of work with different claims, a
Type-II profile the programme's machinery cannot see, and a wall --- Wall~9 ---
forbidding the import of any of it. That chapter ends with a survey. This one
is what happened when the programme went back and pushed on the survey's
conclusions, one at a time, over five consecutive sessions.

The short form is easy to state and worth stating before any of the
mathematics: \emph{nothing moved}. Five sessions, four of them commissioned to
find a way through or around the central obstruction, produced four correct
results and one literature sweep, and not one of them reclassified a single
wall. That is the chapter's subject.

\section{The thesis, stated first}
\label{sec:pt-thesis}

The programme's central obstruction is Wall~1 (\S\ref{sec:wall1},
Chapter~\ref{ch:cz-log}): the sup-norm of a Calder\'on--Zygmund image costs
\(M\Lfac\) and not \(M\), and every enstrophy budget from \Sn{33} onward
inherits the factor. It was identified as defect D8 in the \Sn{31} audit,
consolidated as a wall in \Sn{47}, and has since been tested by every session
that had a plausible reason to think it could be got round.

\begin{record}
\textbf{The claim this chapter makes.} The \Sn{53}--\Sn{57} arc is the
\textbf{fifth, sixth and seventh independent confirmation} --- after \Sn{48},
\Sn{50} and \Sn{52} --- that the obstruction is real and does not move:

\smallskip
\emph{not under reformulation} (\Sn{53}: adapting the rescaling to reach
Type-II does not adapt the tools, it deletes the term they are built on);

\emph{not under the one concrete external candidate profile} (\Sn{54},
\Sn{55}: the only explicit Type-II profile in existence is excluded in the
unforced axisymmetric class, and the forcing it needs instead is bounded below
at full leading order);

\emph{not under the best external methodology this programme could find}
(\Sn{56}, \Sn{57}: located, read in full from primary sources, priced, and
refuted as a route around the wall, with one genuine technique kept).

\smallskip
This is not a breakthrough. It is the most thorough stress-test the
obstruction has received, and it survived all of it.
\end{record}

Two things have to be said immediately, because the claim is easy to inflate.

\textbf{What ``confirmation'' means here.} It does not mean that anything was
proved \emph{about} Wall~1. Wall~1 remains \OPEN{} and structural; no proof in
the document says the logarithm cannot be removed, and none of these five
sessions produced one. What each session did was take a specific candidate
route --- a reformulation, a profile, a methodology --- and establish that it
does not deliver what the wall requires. Each closes with the same ledger line,
in its own words: \emph{Wall~1, Wall~3, Wall~9 unchanged in classification}.
Seven failed routes are not a theorem. They are evidence about where to look,
which is what a walls document is for.

\textbf{Three of the five sessions are not literally about Wall~1.} \Sn{53}
sharpens Wall~3 (the ceiling), \Sn{54} and \Sn{55} sharpen Wall~9
(forcing).\footnote{The counting differs between two of this chapter's own
sources, and the difference is bookkeeping, not substance. The
\texttt{PROGRESS.md} \textsc{key finding} block for \Sn{57} calls that session
``the fifth independent confirmation (after \Sn{48}, \Sn{50}, \Sn{52},
\Sn{53})'' --- four priors, the arc contributing one. The brief commissioning
this chapter counts the wall's own original identification (D8 at \Sn{31},
consolidated at \Sn{47}) as the first, making \Sn{48}, \Sn{50} and \Sn{52} the
second through fourth and the arc the fifth through seventh, with \Sn{54} and
\Sn{55} read as one confirmation and \Sn{56}--\Sn{57} as another. Nothing
mathematical depends on which convention is used; this chapter adopts the
brief's, and records the discrepancy rather than silently picking one.} They
belong in the same arc anyway, for a reason \Sn{53} makes explicit: the walls
are not independent. The machinery of \S33--\S44 is parabolic and Type-I-rate
essential, so a route that escapes Wall~1 by leaving the Type-I class hands the
problem to Wall~3, and a route that escapes Wall~3 by importing a constructed
singularity hands it to Wall~9. The arc tested all three exits, and found each
one closed by a different mechanism.

\begin{figure}[ht]
\centering
\begin{tikzpicture}[
  >={Latex[length=2mm]},
  ev/.style={draw, rounded corners=2pt, font=\scriptsize, align=left,
             inner sep=4pt, text width=64mm},
  neg/.style={ev, draw=impossiblered, thick},
  neu/.style={ev, draw=rulegrey},
  keep/.style={draw, rounded corners=2pt, font=\scriptsize, align=left,
               inner sep=4pt, text width=42mm, draw=provedgreen, very thick},
  ar/.style={->, thick, rulegrey}
]
\node[neu] (s53) at (0,0)
  {\Sn{53} \textbf{Can the machinery be adapted to Type-II?}\\
   \texttt{lem:wpt1-nueff}: any rescaling normalising a supercritical rate
   has \(\nu_{\mathrm{eff}}=\nu A/L\to0\). The limit solves \emph{Euler}};
\node[neg, below=6mm of s53] (s54)
  {\Sn{54} \textbf{Does the one explicit Type-II profile survive unforced?}\\
   \texttt{prop:wpt2-swirl-exclusion}: no. Its swirl
   \(\Gamma\asymp\tau^{-h}\to\infty\) contradicts the swirl maximum principle};
\node[neg, below=6mm of s54] (s55)
  {\Sn{55} \textbf{How much forcing does it actually need?}\\
   \texttt{prop:wpt3-floor}: \(\abs{f_\vartheta}\gtrsim\tau^{-(A+1)}\) at the
   swirl maximum --- full leading order. The source's own exponent agrees};
\node[neu, below=6mm of s55] (s56)
  {\Sn{56} \textbf{Is there anything better outside?}\\
   One candidate in the entire 2024--2026 sweep: a log-weighted BMO depletion
   mechanism, in this document's own direction equation};
\node[neg, below=6mm of s56] (s57)
  {\Sn{57} \textbf{Is it cheaper than what we already have?}\\
   \texttt{thm:s57-independence}: not a relabelling, not comparable, not
   cheaper. Three explicit counterexamples};

\foreach \a/\b in {s53/s54, s54/s55, s55/s56, s56/s57}{\draw[ar] (\a) -- (\b);}

\node[keep, right=7mm of s53] (k53)
  {\textbf{kept:} the trichotomy, and the correct reason the vacuity holds};
\node[keep, right=7mm of s54] (k54)
  {\textbf{kept:} \texttt{cor:wpt2-radius-cap} --- check a candidate profile
   against it \emph{first}};
\node[keep, right=7mm of s55] (k55)
  {\textbf{kept:} the forced swirl maximum principle, a genuine
   generalisation};
\node[keep, right=7mm of s56] (k56)
  {\textbf{kept:} a map of what does \emph{not} exist in the literature};
\node[keep, right=7mm of s57] (k57)
  {\textbf{kept:} the commutator/CRW estimate, verified correct};

\foreach \a/\b in {s53/k53, s54/k54, s55/k55, s56/k56, s57/k57}
  {\draw[->, thick, provedgreen] (\a.east) -- (\b.west);}

\node[draw=impossiblered, very thick, rounded corners=2pt, font=\scriptsize,
      align=center, inner sep=5pt, text width=100mm, below=7mm of s57.south west,
      anchor=north west] (led)
  {\textbf{Ledger after five sessions:} Wall~1 \OPEN{}, Wall~3 \OPEN{},
   Wall~9 \IMPOSSIBLE{} as an import route --- \emph{every classification
   exactly as \Sn{47} left it}};
\draw[ar] (s57) -- (led);
\end{tikzpicture}
\caption[The five-session arc]{The arc, as a chain. Each session asked one
question about a route past the programme's obstructions; each returned a
correct negative, and each left something behind that is worth keeping
independently of the negative. The bottom box is what did not change.}
\label{fig:pt-arc}
\end{figure}

\section{\texorpdfstring{\Sn{53}: what a rescaling costs}{S53: what a rescaling costs}}
\label{sec:pt-s53}

Chapter~\ref{ch:openai}, \S\ref{sec:s49-two-strands}, recorded the fact that
cuts hardest against this programme: the one explicitly constructed singular
profile in the literature is Type-II, and the machinery of \S33--\S39 is blind
to Type-II by construction. The document's own statement of why is a single
sentence in \texttt{rem:s39-distance}: ``a Type-II singularity would rescale to
nothing''. \Sn{53} was commissioned to find out whether that sentence is true,
and whether the blindness is fundamental or an artefact.

It is neither, as posed, and the true statement is \emph{worse} for the
programme than the one it replaces.

\subsection{The price of a rescaling, computed}

The escape everyone reaches for is to rescale differently. Give the rescaling
two parameters instead of one --- an amplitude \(A\) and a length \(L\) --- and
ask what equation the limit satisfies. Writing
\(v(y,s)=A\,u(x^*+Ly,\,T+\tau s)\) and normalising the nonlinearity by
\(\tau=AL\), \texttt{lem:wpt1-nueff} gives the answer in one line:
\[
  \nu_{\mathrm{eff}} \;=\; \nu\,\frac{A}{L}.
\]
The programme's standard rescaling is the case \(A=L\), which returns
\(\nu_{\mathrm{eff}}=\nu\). Every other member of the family changes the
viscosity, and the direction of the change is forced by the blow-up rate it is
asked to normalise. Requiring \(\sup\abs{u}\asymp(T-t)^{-\beta}\) to rescale to
size one gives \(A=L^{\beta/(1-\beta)}\), hence
\(\nu_{\mathrm{eff}}/\nu=L^{(2\beta-1)/(1-\beta)}\), and
\texttt{cor:wpt1-trichotomy} reads it off.

\begin{table}[ht]
\centering
\small
\begin{tabular}{@{}l l p{0.44\textwidth}@{}}
\toprule
Rate \(\beta\) & \(\nu_{\mathrm{eff}}\) as \(L\downarrow0\) & What the limit
  object solves \\
\midrule
\(\beta<\tfrac12\) & \(\to\infty\)
  & Stokes / heat. No blow-up mechanism survives \\[3pt]
\(\beta=\tfrac12\) & \(=\nu\)
  & \textbf{Navier--Stokes.} \(A=L\); the \emph{unique} member of the family
    that preserves the equation, and exactly the programme's own rescaling \\[3pt]
\(\tfrac12<\beta<1\) & \(\to0\)
  & \textbf{Euler.} This is Seregin's zoom, with \(\beta=\alpha/(\alpha+1)\) \\[3pt]
\(\beta\geq1\) & ---
  & \textbf{No member of the family works.} At \(\beta=1\) normalisation forces
    \(L=1\), i.e.\ no spatial zoom; for \(\beta>1\) it forces
    \(\tau=AL\to\infty\), contradicting \(\tau\leq T\) \\
\bottomrule
\end{tabular}
\caption[The rescaling trichotomy]{\texttt{cor:wpt1-trichotomy}. The
programme's rescaling is not one choice among many: it is the only one that
leaves the equation alone. The last row is a scope limit on the Euler-zoom
method itself and is stated in no source \Sn{53} read.}
\label{tab:pt-trichotomy}
\end{table}

Three independent checks were run and all three agree: \(A=L=\lambda\) returns
the document's own rescaling; Seregin's Euler scaling \(A=\lambda^\alpha\),
\(L=\lambda\) returns \(\nu_{\mathrm{eff}}=\nu\lambda^{\alpha-1}\to0\) for
\(\alpha>1\); and his \(f\)-form returns \(\nu_{\mathrm{eff}}=\nu f(\lambda)\).

\subsection{The correction, reported as the session was asked to report it}

\begin{record}
\textbf{\texttt{rem:s39-distance}(1) is right about the conclusion and
imprecise about the reason,} and the correct statement --- \Sn{53}'s
\texttt{rem:wpt1-correction}, which future sessions should quote instead --- is
this:

\smallskip
\emph{A rescaling that normalises a supercritical blow-up rate in \(L^\infty\)
necessarily has \(\nu_{\mathrm{eff}}=\nu A/L\to0\), so its limit object solves
\textbf{Euler}. The vacuity is not that nothing survives; it is that what
survives is \emph{inviscid} --- and every tool of \S33--\S44 is parabolic:
Moser iteration, the magnitude Caccioppoli inequality, De Giorgi level
iteration, \texttt{lem:local-enstrophy-typeI}, CKN \(\varepsilon\)-regularity,
backward uniqueness. None of them says anything about a solution of the Euler
equations.}
\end{record}

That is a strictly stronger obstruction than ``there is no limit object''. A
limit object that exists and lies outside every estimate the programme owns is
worse than no limit object at all, because it removes the hope that a
compactness argument could be repaired. \Sn{53} also found a second respect in
which the old sentence was too strong: under the weaker Caffarelli--Kohn--%
Nirenberg hypothesis \(g<\infty\), the \emph{same} \(\lambda\)-family is bounded
in \(L^\infty_tL^2_{\mathrm{loc}}\cap L^2_t\dot H^1_{\mathrm{loc}}\) and the
classical compactness does produce an ancient suitable weak Navier--Stokes
solution, viscosity retained.

\subsection{The loophole, and why it closes on dimensional grounds}

If the Euler zoom must throw the viscosity away, the natural hope is to find one
that still carries something this programme controls. In the axisymmetric class
the obvious candidate is the swirl \(\Gamma=r\,u_\vartheta\), which is free by
the maximum principle. It does not work, and \texttt{cor:wpt1-swirl} gives the
reason in one line: under the two-parameter family,
\[
  \Gamma_{\mathrm{new}} \;=\; \frac{A}{L}\,\Gamma
  \;=\; \frac{\nu_{\mathrm{eff}}}{\nu}\,\Gamma ,
\]
because \(\Gamma\) and \(\nu\) have the same dimension \(L^2/T\), so
\(\Gamma/\nu\) is invariant across the whole family. Hence any zoom that keeps
the swirl keeps the viscosity --- that is, is the programme's own rescaling, and
needs Type-I. ``An Euler zoom that keeps the swirl'' is dimensionally
impossible, not merely unknown.\footnote{The conclusion in the direction
``\(\nu_{\mathrm{eff}}\to0\Rightarrow\Gamma\to0\)'' carries a hypothesis that
the \Sn{53} log states without: the maximum principle \texttt{lem:s51-gamma}
holds for the \emph{unforced} equations. \Sn{54} found a concrete witness that
the hypothesis is load-bearing, and the proof document's version of the
corollary cites the lemma correctly. The document governs; the \Sn{53} log's
phrasing should not be reused. Note also that \texttt{lem:s51-gamma} and its
neighbours carry the label prefix \texttt{s51} while the session that wrote
them was \Sn{52} --- a pre-existing off-by-one between the document's headings
and the session logs, recorded here so it is not rediscovered.}

\subsection{The verdict, and the region nobody owns}

\Sn{53}'s verdict was referral, not a target: the branch of mathematics that
owns Type-II for true three-dimensional Navier--Stokes is \emph{Liouville
theorems for ancient solutions of the Euler equations in weighted classes
defined by a blow-up scenario}, and it is one author's programme --- four
papers by Seregin, all of them \emph{scenario-conditional}. No unconditional
exclusion of Type-II exists in that literature, and none is claimed there. This
programme has no foothold in it: it has never written an Euler estimate.

One structural finding deserves to be carried, because it means the two
literatures do not tile the space they appear to. The Type-II results are
stated in a \emph{third} convention --- Caffarelli--Kohn--Nirenberg
scale-invariant energies --- and temporal Type-I implies the CKN version, so
Seregin's Type-II class is strictly \emph{smaller} than this document's. Between
them lies a region, \(\mathrm{(R2)}\): temporally Type-II but CKN Type-I,
addressed by neither. Whether it is populated is unknown. Any future claim to
have ``covered Type-II'' has to say what it does about it.

And \Sn{53} kept ``open'' and ``expected'' apart, deliberately. Nothing read
excludes Type-II; equally, no source read asserts that it occurs. The strongest
unforced numerical candidate in the literature is Type-I-shaped, and the only
Type-II profile is forced.

\section{\texorpdfstring{\Sn{54}: the one concrete candidate, and what killed
it}{S54: the one concrete candidate, and what killed it}}
\label{sec:pt-s54}

\Sn{53} left one obvious next question. There exists, in the world, exactly one
explicit Type-II profile for a Navier--Stokes-type system: the collapsing
anisotropic vortex core of Chapter~\ref{ch:openai}, \S\ref{sec:their-mechanism}.
Does it survive as an ancient solution of the unforced equations?

The first thing \Sn{54} did was fetch the primary document rather than work from
the programme's own transcription of it. The transcription --- carried since
\Sn{40} from a screenshot --- turned out to be \emph{correct in every symbol},
and the governing profile equations turned out to be stated explicitly in the
source rather than left implicit in its asymptotics. One fact was missing from
the programme's record, and it governed everything downstream: the leading
balance \textbf{retains radial viscosity and omits axial viscosity}. The source
says so three separate ways, including the dimensional statement
\(\mathrm{Re}_r=O(1)\) already reproduced in \S\ref{sec:their-mechanism}.

\subsection{The rescaling with two lengths, and a limit that is not Euler}

\texttt{lem:wpt2-aniso} generalises \texttt{lem:wpt1-nueff} to independent
radial and axial scales. With \(r=L_r\rho\), \(z=L_z\zeta\), and amplitudes
\(A_r,A_\vartheta,A_z\), the limit is incompressible \emph{if and only if}
\(A_rL_r=A_zL_z=:\Lambda\), and normalising the advection by \(\Lambda\)
produces not one effective viscosity but two:
\[
  \nu_r=\nu\,\frac{A_r}{L_r},\qquad \nu_z=\nu\,\frac{A_z}{L_z}.
\]
\texttt{lem:wpt1-nueff} is the isotropic special case, and was re-derived
independently before being used.

Fed the profile's own exponents, \texttt{prop:wpt2-limit} returns three things,
all of which came out against the session's expectation. The incompressibility
constraint holds \emph{automatically}, because the profile's amplitude
anisotropy is exactly conjugate to its length anisotropy --- \(A+D=1\), which is
the source's own identity. The radial viscosity is preserved \emph{exactly},
\(\nu_r=\nu\), while \(\nu_z=\nu\tau^{2h}\to0\). And the radial momentum
equation degenerates, its pressure and centrifugal coefficients both diverging,
leaving the cyclostrophic balance \(\partial_\rho\tilde q=v_\vartheta^2/\rho\)
--- which is again the source's own identity.

So the limit exists, is closed and nontrivial, and \textbf{does not solve
Euler}. Structurally it is a boundary-layer (Prandtl-type) reduction: momentum
degenerated to a constraint in the thin direction, thin-direction velocity
recovered from continuity, diffusion in the thin direction only. Two cautions
were recorded with it, and both belong here. No source read writes this system,
so the identification is structural and is offered as such. And well-%
\emph{defined} is not well-\emph{posed}: Prandtl's system is linearly ill-posed
in Sobolev spaces, by Gérard-Varet--Dormy, whose proof exploits exactly the
features this limit has. That result is not transported, and no claim is made
that it applies.

The consequence for the Liouville programme is negative in both directions.
Those theorems need an \emph{isotropic} Euler zoom, and one length cannot match
two: scaling by \(\ell_r\) gives a steady, \(\zeta\)-independent, infinite-%
energy columnar Euler flow lying in no weighted class any of the four theorems
uses; scaling by \(\ell_z\) collapses the core onto the axis and the limit is
trivial, so nontriviality --- the hypothesis every Liouville argument
contradicts --- fails outright.

\subsection{The exclusion, which needed none of that}

The profile is excluded, and the argument uses none of the preceding
subsection. It is three lines of the swirl maximum principle.

\begin{record}
\textbf{\texttt{prop:wpt2-swirl-exclusion}.} The profile's angular momentum is
\(\Gamma^{(0)}=r\,u_\vartheta^{(0)}=q^{-h}H(X,\eta)\) --- an identity stated
verbatim in the source's own proof, and re-derived independently here --- so
\(\sup\abs{\Gamma^{(0)}}\gtrsim\tau^{-h}\to\infty\). But for smooth
axisymmetric solutions of the \emph{unforced} equations,
\texttt{lem:s51-gamma} gives \(\norm{\Gamma(t)}_\infty\leq
\norm{\Gamma(0)}_\infty\), and Schwartz data have
\(\norm{\,\abs{x}u_0}_\infty<\infty\). \textbf{Contradiction.} The profile
cannot be the blow-up asymptotics of an unforced axisymmetric Navier--Stokes
solution.
\end{record}

The geometric form, \texttt{cor:wpt2-radius-cap}, is the part worth carrying
forward. The maximum principle gives \(\abs{u_\vartheta}\leq\Gamma_0/r\) ---
that is, the swirl component automatically satisfies the \emph{spatial} Type-I
bound. So on a core of radius \(\ell_r\) the swirl is capped at
\(\Gamma_0/\ell_r\), and an unforced axisymmetric core with the Type-I radial
geometry \(\ell_r\asymp\tau^{1/2}\) can carry only Type-I swirl. Type-II swirl
requires a strictly thinner core, \(\ell_r\lesssim\tau^{1/2+h}\). The profile
has \(\ell_r\asymp\tau^{1/2}\) together with
\(\abs{u_\vartheta}\asymp\tau^{-1/2-h}\), and misses the cap by exactly the
factor \(\tau^{-h}\) that makes it Type-II.

\begin{figure}[ht]
\centering
\begin{tikzpicture}[>={Latex[length=2mm]}, scale=1.0]
  \draw[->, thick, gray] (0,0) -- (8.7,0)
    node[right, font=\scriptsize] {\(\log\tfrac1\tau\)};
  \draw[->, thick, gray] (0,0) -- (0,4.5)
    node[above, font=\scriptsize] {\(\log\sup\abs{u_\vartheta}\)};

  % shaded wedge between the two rates
  \fill[impossiblered, opacity=0.10]
    (0,0.35) -- (7.8,4.15) -- (7.8,3.20) -- cycle;

  % Type-I cap: slope 1/2
  \draw[provedgreen, very thick] (0,0.35) -- (7.8,3.20);
  % Type-II profile: slope 1/2+h
  \draw[impossiblered, very thick] (0,0.35) -- (7.8,4.15);

  \draw[<->, thick] (8.05,3.20) -- node[right, font=\scriptsize]
    {\(\tau^{-h}\)} (8.05,4.15);

  \node[impossiblered, font=\scriptsize, anchor=south west, align=left]
    (plab) at (0.9,3.05) {the profile: \(\tau^{-1/2-h}\)};
  \draw[impossiblered, thin] (plab.south east) -- (3.55,2.15);
  \node[provedgreen, font=\scriptsize, anchor=north west, align=left]
    at (3.4,1.35) {cap from the maximum principle:\\
      \(\abs{u_\vartheta}\leq\Gamma_0/\ell_r\asymp\tau^{-1/2}\)};
\end{tikzpicture}
\caption[The swirl radius cap]{\texttt{cor:wpt2-radius-cap} in one picture,
both lines drawn at the same core radius \(\ell_r\asymp\tau^{1/2}\). In the
unforced axisymmetric class the swirl maximum principle caps
\(\abs{u_\vartheta}\) by \(\Gamma_0/r\), so a core of Type-I radius can carry
only Type-I swirl. The shaded wedge is the excess \(\tau^{-h}\) that makes the
candidate profile Type-II --- and it is exactly the amount by which the profile
violates a free conservation law.}
\label{fig:pt-radius-cap}
\end{figure}

\subsection{Four things this is not}

\Sn{54}'s own statement of its scope is reproduced here without softening,
because the result is narrow and the temptation to quote it broadly is real.

\begin{enumerate}[label=\textup{(\alph*)},leftmargin=2.4em]
\item \emph{Axisymmetric class only.} \texttt{lem:s51-gamma} is. Nothing is said
  about a non-axisymmetric flow with the same scales.
\item \emph{Unforced only} --- which is the point of the exercise, but which
  also means \textbf{this is not a statement about the source's theorem}. That
  flow is forced \emph{and} non-axisymmetric, so the maximum principle never
  applied to it. No contradiction with it is asserted or implied. Wall~9 governs
  (\S\ref{sec:wall9}).
\item \emph{One profile family, not Type-II.} Any axisymmetric core with
  \(\ell_r\lesssim\tau^{1/2+h}\) escapes the cap entirely. Wall~1, Wall~3 and
  region \(\mathrm{(R2)}\) are untouched.
\item \emph{Elementary.} The swirl maximum principle is classical and was
  already in the document. The content is the check, not the machinery, and no
  credit is claimed for the latter.
\end{enumerate}

The transferable lesson is a piece of process, and it is the one thing from
\Sn{54} a future session should actually reuse: the Liouville programme was the
wrong tool, two independent computations say so, and what settled the question
was a conservation law the class supplies for free. \emph{Before importing
profile-extraction machinery against a candidate profile, check it against
\texttt{cor:wpt2-radius-cap} first.}

\section{\texorpdfstring{\Sn{55}: putting a number on Wall 9}{S55: putting a number on Wall 9}}
\label{sec:pt-s55}

Wall~9 says the forcing in such a construction is ``load-bearing, not
decorative'' (\S\ref{sec:wall9}). That is a qualitative statement, inherited
from \Sn{40}'s reading, and a qualitative statement is exactly the kind of thing
that decays into a slogan. \Sn{55} turned it into a number.

\subsection{The forced swirl maximum principle}

The tool is a generalisation of \texttt{lem:s51-gamma} that did not exist in the
document. \texttt{lem:wpt3-forced-swirl}: for axisymmetric flow, adding a force
\(f\) adds \textbf{exactly} \(r f_\vartheta\) to the swirl equation and nothing
else. Neither \(f_r\) nor \(f_z\) appears, and there is still no pressure term.
The three places this could fail --- pressure, advection, the vector Laplacian
--- were each checked rather than assumed.

\texttt{lem:wpt3-forced-max} then runs the maximum principle with the force in
place, on three hypotheses and no more: smoothness; attainment of the maximum
with \emph{uniformly} compact maximiser sets; and positivity of the running
maximum. Two features are worth recording. The maximiser is automatically off
the axis, since \(\Gamma=r u_\vartheta=O(r^2)\) there --- so the singular
coefficient \(-2\nu/r\) is never evaluated at \(r=0\), and the
\(\varepsilon\)-exhaustion used in the unforced lemma's own proof is not needed.
The forced statement is, in that one respect, \emph{cleaner} than the statement
it generalises. And the derivative of the maximum is handled by Danskin, both
one-sided derivatives written out, with an integrated form that assumes no
differentiability anywhere; at \(f\equiv0\) it returns \texttt{lem:s51-gamma}
exactly, so it is a genuine generalisation rather than a parallel statement.

\begin{record}
\textbf{The commissioning brief's own statement of the target was false, and
had to be corrected.} It asked for
\(\tfrac{\dee}{\dee t}\norm{\Gamma}_\infty\leq(rf_\vartheta)(x_m(t),t)\). For
the absolute-value norm that is wrong: if the norm is attained where
\(\Gamma<0\), that point is a global \emph{minimum}, the diffusion term there
is \(\geq0\), and the inequality \emph{reverses}. The correct statements are the
brief's inequality for the \emph{signed} maximum, and a version carrying a sign
factor for the norm. Nothing downstream changed --- the floor is stated for the
signed maximum --- but the erroneous form should not be propagated. Recorded in
the document as \texttt{rem:wpt3-handover-sign}.
\end{record}

\subsection{The floor}

Because the leading field is exactly self-similar, its swirl maximum is not
merely comparable to a power of \(\tau=1-t\) but \emph{equal} to one:
\(M(t)=\sup_x\Gamma^{(0)}=c_*\tau^{-h}\) exactly, with \(c_*\) attained in the
interior of the profile domain. That exactness is what makes the result a floor
rather than a scaling heuristic. Differentiating and applying the forced maximum
principle gives \texttt{prop:wpt3-floor}:
\[
  \abs{r f_\vartheta^{(0)}} \;\geq\; h\,c_*\,\tau^{-(A+1/2)},
  \qquad
  \abs{f_\vartheta^{(0)}} \;\geq\; \frac{h\,c_*}{c_r}\,\tau^{-(A+1)}
  \;=\;\frac{h\,c_*}{c_r}\,\tau^{-(3/2+h)},
\]
pointwise at the running swirl maximum, whose radius \(c_r\tau^{1/2}\) is an
\emph{output} of the computation and not an assumption fed into it.

The calibration is what gives the exponent its meaning. Every individual term
of the leading field's own azimuthal momentum equation --- the time derivative,
the radial advection, the radial viscous term --- is of size
\(\tau^{-(A+1)}\), while the axial viscous term is smaller by \(\tau^{2h}\). So
the floor says the azimuthal residual is of \emph{full leading order in its own
momentum equation}: not a perturbation, not lower order, and not made so by any
choice of the free profile data. That is the quantitative content of Wall~9's
``load-bearing''. As a by-product the same computation independently reproduces
the source's own structural claim that radial viscosity is retained in the
leading balance and axial viscosity is not.

\subsection{The comparison, and what it does not say}

The floor was derived from the swirl maximum principle and a single exponent.
The source states its own forcing scaling, in a different place and for a
different reason, as a stress exponent \(T=q^{-A-1/2}T_0\). Converting gives
\(\abs{f_{\vartheta,B}}\asymp\tau^{-(A+1)}\) for the axisymmetric background:
\textbf{the same two exponents, by a genuinely different route}, with radial
integration of the residual on one side and a maximum principle on the other.
Neither derivation uses the other's input.

Four qualifications travel with that, all of them from the session's own record.
The comparison is \textbf{exponent-level only}: the floor's constant carries
\(h<1/100\) and the source quotes no numerical value, so constants are not
compared. The statement is \textbf{pointwise, at one moving point} --- a set of
measure zero, from which no positive-measure and no \(L^q\) statement for
\(q<\infty\) follows, a spike of measure \(\varepsilon\) tracking the maximiser
being enough to satisfy it. There is \textbf{no tension with the source's main
theorem}, whose final force is flat and whose full field is non-axisymmetric, so
the lemma does not apply to it at all; what the floor bounds is the residual of
the \emph{axisymmetric background}, an intermediate object the construction
never claims is small. And \textbf{no error is asserted, and none was found}. A
single derived requirement the source does not state --- that the swirl maximum
is forced into the active annulus where the pulses act,
\texttt{cor:wpt3-annulus} --- is recorded as a fact checkable in one line by
anyone holding the profile, explicitly not as a defect.

\begin{obsv}[The sharpest available form of Wall 9]
\label{obs:pt-wall9-sharp}
Read in the useful direction, the floor explains \emph{why} such constructions
break axisymmetry: any mechanism carrying an axisymmetric background's swirl
from \(O(1)\) to \(\tau^{-h}\) must inject azimuthal momentum at full leading
order, and the only free conservation law of the class forbids doing it for
nothing. Breaking axisymmetry is not a convenience of the construction; it is
what the swirl maximum principle costs.

\smallskip
This is documentation-sharpening, and the session classified it as such. It is
not a step toward alternatives (A) or (B) --- it proves a \emph{lower} bound on
a forcing, which is the opposite direction of travel --- and Wall~9's
classification is unchanged.
\end{obsv}

\section{\texorpdfstring{\Sn{56}: what the literature actually contains}{S56: what the literature actually contains}}
\label{sec:pt-s56}

At this point the internal toolkit was, in the sense \Sn{52} and \Sn{53} used
the word, exhausted. \Sn{56} was a strategic audit: sweep the 2024--2026
literature against the now-precise statement of the obstruction and report what
exists.

It returned one find, one honest negative, and one piece of context that is
uncomfortable and belongs in the record anyway.

\textbf{The find.} Two preprints of a single author propose a conditional route
to global regularity that operates \emph{in this document's own direction
equation}, at \emph{this document's own critical scale}, and attacks the
logarithm at exactly the point Wall~1's consequence clause declares permitted:
the gradient of \(u\) is never recovered by an \(L^\infty\)
Calder\'on--Zygmund bound. The mechanism recasts the stretching eigenvalue
\(\alpha=\ehat\cdot S\ehat\) as a singular-integral commutator and pays for it
with Coifman--Rochberg--Weiss; the hypothesis carrying it is that the vorticity
direction lies in a \emph{logarithmically weighted} space of bounded mean
oscillation, a condition that fails the Dini criterion and is therefore far
weaker than Constantin--Fefferman. The identification of the direction equation
was verified term for term against the document's own \texttt{eq:dir-eq}, and
the concentration regime against the document's own criticality
\(M\,(\rstar)^2=1\). It was the only methodology the sweep found that
plausibly satisfies Wall~1's escape clause.

\textbf{The counterweight, recorded at the same time and treated as
load-bearing:} unrefereed, single author, sweeping conclusion.

\textbf{The negative.} On the branch \Sn{53} named as owning Type-II ---
Liouville theorems for ancient Euler flows --- there is \emph{nothing new
whatsoever}. The systematic sweep returns exactly the four papers \Sn{53} had
already read, and the 2024--2026 ``ancient solutions'' literature is essentially
all geometric flows. The branch this programme was told to refer the question to
has not moved in the interval.

\textbf{Two methodologies closed by citation rather than opinion.} A
shell-model blow-up result whose interactions are, in its own words, organically
represented in the true Euler/Navier--Stokes nonlinearity means no purely
cascade-budget argument can distinguish Navier--Stokes from a blowing-up object.
And computer-assisted proof on the regularity side exists but is structurally a
posteriori and short-time: it can exclude blow-up for fixed data on a fixed
interval and can never deliver a universally quantified statement.

\textbf{One ledger correction, made here and worth flagging because it cuts in
the programme's disfavour.} Wall~9 does \emph{not} govern the separate,
\emph{unforced} Euler construction by the same group; it continues to govern the
forced Navier--Stokes paper of Chapter~\ref{ch:openai} exactly as written. The
two are different papers, and reading the wall as covering both would be an
error in this programme's own favour.

\textbf{And the context.} Every AI-assisted effort the sweep found targets
blow-up; none targets regularity. \Sn{56} records the structural reason rather
than treating it as fashion: blow-up is an existence statement with a
certificate, regularity is a universally quantified statement without one. That
asymmetry will not reverse on its own, and this programme is on the side of it
that has no certificate to offer.

\section{\texorpdfstring{\Sn{57}: pricing the candidate}{S57: pricing the candidate}}
\label{sec:pt-s57}

\Sn{56} ended with a recommendation and a prediction. The recommendation was one
gated session to \emph{price} the candidate --- not to adopt it --- under the
standard \Sn{50} set: a new target must be proved cheaper, not asserted cheaper.
The prediction was that the answer would be a fifth avatar of a requirement the
programme already owns, the pattern of \Sn{50} and \Sn{52}, which would itself
be a legitimate result.

\Sn{57} refuted the prediction and confirmed the conclusion. The hypothesis is
\emph{not} a relabelling of the programme's own \(\sstar\)-decay --- the two are
logically independent --- and it is still not cheaper.

\subsection{The dictionary, and the one implication that survives}

The proposed dictionary was a Poincar\'e inequality: gradients control
oscillation, so a smallness hypothesis on the direction's mean oscillation
should read as a smallness hypothesis on \(\sstar\). Poincar\'e does run in that
direction, and \texttt{lem:s57-poincare} makes the surviving implication exact:
if the \emph{unrestricted} deficit satisfies
\(\tilde\sigma_\rho(y)\leq K\abs{\log\rho}^{-2}\) for \emph{every} ball
\(B_\rho(y)\subseteq B_{\rstar}(x^*)\), then the local log-weighted seminorm of
\(\ehat\) is at most \(\tfrac14K^{1/2}\). At the top scale this is precisely the
sketch \Sn{56} proposed --- but only as a \emph{sufficient} condition, only for
the unrestricted deficit, and only at one ball.

Everything else fails, and the failures are not quantitative near-misses. Three
explicit fields settle it, all smooth, all divergence-free, hence all the
vorticity of a solenoidal velocity field.

\begin{table}[ht]
\centering
\small
\begin{tabular}{@{}l >{\raggedright\arraybackslash}p{0.29\textwidth}
                     >{\raggedright\arraybackslash}p{0.24\textwidth}
                     >{\raggedright\arraybackslash}p{0.28\textwidth}@{}}
\toprule
& Field & What it does & What it refutes \\
\midrule
\textbf{(C1)}
  & \(\omega=M(\cos g_\delta(x_3),\sin g_\delta(x_3),0)\),
    \(g_\delta(s)=\frac{\sin(s/\delta)}{2\abs{\log\delta}}\)
  & seminorm \(\leq2\) uniformly, while \(\sstar\to\infty\)
  & the hypothesis carries \emph{no} \(L^2\) information about
    \(\nabla\ehat\) at all \\[10pt]
\textbf{(C2)}
  & \(\omega=M\cos(x_3/\delta)\,e_1\)
  & \(\sstar=0\) exactly, while the seminorm is \(+\infty\)
  & the level-set indicator of \(\sstar\) is load-bearing, not cosmetic \\[10pt]
\textbf{(C3)}
  & \(\omega=e_3+\nabla\times(\psi_\delta e_1)\),
    \(\psi_\delta=\delta\chi_0(\abs{x-y_0}/\delta)\)
  & \(\sstar\to0\), while the hypothesis fails at scale \(\delta\) by a factor
    \(\abs{\log\delta}\)
  & one scale does not give the scale family \\
\bottomrule
\end{tabular}
\caption[The three counterexamples]{The three fields of
\texttt{prop:s57-no-gradient}, \texttt{prop:s57-indicator} and
\texttt{prop:s57-onescale}, which together prove
\texttt{thm:s57-independence}: neither condition implies the other. In (C1) and
(C3) the indicator is inert by construction, so the failure cannot be blamed on
it; (C2) is the case where it does all the work.}
\label{tab:pt-counterexamples}
\end{table}

\begin{obsv}[What (C2) says about the programme's own definition]
\label{obs:pt-indicator}
The restriction to \(\{\abs\omega\geq M/4\}\) in the definition of \(\sstar\)
has always looked like a convenience. It is not. Near any \emph{transversal
zero} of \(\omega\) the direction has an order-one jump at every scale, so the
unrestricted deficit is \textbf{generically infinite} --- and the indicator is
what keeps \(\sstar\) finite at all. Any unrestricted-oscillation hypothesis
imported into this programme therefore needs a zero-free core supplied as a
separate assumption. That is a fact about the programme's own machinery,
discovered by pricing somebody else's.
\end{obsv}

The honest translation, once the three failures are accounted for, is that the
external hypothesis differs from \(\sstar\)-decay in three independent ways:
weaker in norm (\(L^1\) oscillation against an \(L^2\) gradient bound --- the
one genuine relaxation, and the reason it carries no gradient information),
stronger in scale (a family on \((0,\rstar]\), not a statement at \(\rstar\)),
and stronger in support (unrestricted, where \(\sstar\) is level-restricted).
Incomparable, in short. And against the wider family, the verdict is the one the
programme has now reached four times: a scale-invariant no-concentration
condition, imposed rather than proved, not established for Navier--Stokes
solutions --- by the companion paper's own admission --- and not shown cheaper
than what it would replace. What is new is only that this one is not a
relabelling.

\subsection{The byproduct that is worth keeping}

The commutator step was checked in isolation, link by link, and \textbf{it
holds}. The cancellation is Constantin--Fefferman's, correctly restated in
commutator dress; the smallness of the symbol uses the weight's monotonicity in
the right direction; the extension step is legitimate on a ball; and the
Coifman--Rochberg--Weiss application reduces, via the operator identity
\([R_iR_j,b]=R_i[R_j,b]+[R_i,b]R_j\), to the classical bound for Riesz
transforms, so no general Calder\'on--Zygmund commutator theory is needed. One
cosmetic error in the source was noted and does not affect the conclusion.

\begin{record}
\textbf{What the programme keeps from \Sn{57}, independent of the paper's
overall fate.} \emph{The stretching eigenvalue is a commutator whose symbol is
the direction field, so any smallness of the direction's \textup{BMO} norm at
small scales is a linear gain on \(\alpha\) in a restricted critical Lorentz
norm.}

\smallskip
This is the first estimate on record in this programme that extracts a gain
from the direction field \emph{without} an \(L^\infty\) Calder\'on--Zygmund
bound --- that is, inside the class Wall~1's own consequence clause declares
permitted.

\smallskip
\textbf{And it is not by itself a route around Wall~1.} The smallness it
consumes is not available from \(\sstar\) (Table~\ref{tab:pt-counterexamples}),
and it is not available from the companion paper either, whose transfer theorem
is conditional on four assumed hypotheses and whose own final section exhibits
two natural critical cores --- a point-concentrated tube, a Beltrami core ---
where those hypotheses fail. The propagation question is open, and the source
says so itself.
\end{record}

\subsection{Two corrections \texorpdfstring{\Sn{57}}{S57} made to \texorpdfstring{\Sn{56}}{S56}}

Both were found by going to primary sources rather than to the previous
session's rendering of them, which is the discipline that produced them.

First, \Sn{56} rendered the central theorem's hypothesis as a Lorentz-space
condition on \(\omega\). It is not: it is a structural definition that
additionally supplies a \emph{containment} of the super-level sets in a ball of
a specific radius --- information a norm bound does not give, since a norm
bounds the measure and not the location, and information the source's own
argument needs.

Second, and more to the point, \Sn{56}'s load-bearing counterweight was pointed
at the wrong half of its citation. The paper \Sn{56} invoked as having
undermined the same author's earlier claims has two halves: its negative verdict
concerns certain \emph{a priori} sparseness classes specifically, while its
first half \emph{positively reproves}, as an explicit alternative to the
author's own approach, the conditional criterion that the new preprint's endgame
actually uses. The prior remains a fair caution about interpretation. It does
not attach where \Sn{56} attached it.

\begin{record}
\textbf{The source hierarchy, applied.} Where \Sn{56} and \Sn{57} disagree,
\Sn{57} governs --- later session, primary sources, and the corrections are in
the direction of \emph{weakening} the programme's own counterweight against an
external claim, not strengthening it. Where a session log and the proof document
disagree, the document governs (\S\ref{sec:source-hierarchy}). Both rules were
needed in this chapter: once for \Sn{53}'s statement of the swirl corollary, and
once here.
\end{record}

\section{What the arc changed, and what it did not}
\label{sec:pt-ledger}

\begin{table}[ht]
\centering
\small
\begin{tabular}{@{}l p{0.34\textwidth} p{0.34\textwidth}@{}}
\toprule
Session & What it settled & What it left standing \\
\midrule
\Sn{53}
  & The vacuity against Type-II is derived, not asserted, and its reason is
    that the limit object is \emph{inviscid}
  & Wall~3 \OPEN{}. Type-II unreachable; region \(\mathrm{(R2)}\) owned by
    nobody \\[3pt]
\Sn{54}
  & The one explicit Type-II profile is excluded in the unforced axisymmetric
    class, by a free conservation law
  & Wall~1, Wall~3 untouched. Not a statement about the source's theorem \\[3pt]
\Sn{55}
  & Wall~9's ``load-bearing'' is a number: \(\tau^{-(A+1)}\), full leading
    order, matched exactly by the source's own exponent
  & Wall~9 \IMPOSSIBLE{} as an import route, classification unchanged \\[3pt]
\Sn{56}
  & The external landscape, mapped: one candidate, nothing at all on the
    ancient-Euler branch
  & Internal toolkit exhausted; no wall touched \\[3pt]
\Sn{57}
  & The candidate is logically independent of \(\sstar\)-decay and not cheaper;
    its commutator step is correct and is kept
  & Wall~1 \OPEN{} and structural, exactly as \Sn{47} classified it \\
\bottomrule
\end{tabular}
\caption[The arc's ledger]{Five sessions, five correct results, no
reclassification. The right-hand column is the one the chapter is about.}
\label{tab:pt-ledger}
\end{table}

Two amendments to earlier chapters follow from the arc and are made there rather
than here: the statement in \S\ref{sec:wall1} that a symmetry-reduced
reformulation is ``untested'' is no longer true, and the walls table of
\S\ref{sec:walls-table} has since gained rows this book does not reproduce.
Both are noted in place.

\begin{obsv}[Why a negative arc is worth a chapter]
\label{obs:pt-why}
A reader is entitled to ask what five sessions of confirmed negatives are doing
in a book, and the answer is the same one Chapter~\ref{ch:walls} gives for the
walls document itself: a route that is known to be closed costs nothing to
avoid, and a route that is merely \emph{suspected} to be closed costs a session
every time somebody has the idea again.

\smallskip
There is a sharper reason here. Before this arc, the programme's position on
Type-II rested on one sentence of its own, unverified --- ``a Type-II
singularity would rescale to nothing'' --- and its position on the external
landscape rested on a reading of one expository document. Both have now been
checked against primary sources, and both turned out to be \emph{right in
conclusion and wrong in reason}. That is the characteristic outcome of this
arc, and it is a better thing to know than a slogan that happens to be true.

\smallskip
What the arc did \emph{not} produce is any movement toward alternatives (A) or
(B). The distance recorded in Chapter~\ref{ch:where-it-stands},
\S\ref{sec:ceiling-vs-prize}, is unchanged, as it has been since \Sn{39}. Five
more sessions of correct mathematics did not shorten it by anything at all, and
the honest summary of the arc is the one its own logs use: we tested this as
hard as we know how, and the wall held.
\end{obsv}
